\chapter{Metodologia e Arquitetura Proposta}
\label{arquitetura}

    Este capítulo detalha a arquitetura do sistema desenvolvido, denominado \textit{raft-ha-manager}, e os métodos empregados para garantir a alta disponibilidade do banco de dados PostgreSQL em um ambiente \verb|SCADA|. Serão apresentadas as características do ambiente, os mecanismos de detecção de falhas, a dinâmica do algoritmo de consenso e a configuração do ambiente de simulação utilizado para validação.

    \section{Visão Geral da Solução}
    
        A solução proposta baseia-se na criação de um agente orquestrador customizado, o \textit{raft-ha-manager}, desenvolvido na linguagem Java. Este agente é responsável por interligar o ciclo de vida do banco de dados PostgreSQL com uma implementação própria do protocolo de consenso Raft, utilizando chamadas de procedimento remoto via \verb|gRPC| para comunicação entre os nós. A escolha do Java como linguagem de implementação decorreu da sua robustez no desenvolvimento de sistemas distribuídos de longa duração, da maturidade do ecossistema de bibliotecas disponíveis e da portabilidade garantida pela \verb|JVM| entre diferentes sistemas operacionais.
    
    \section{Caracterização e Detecção de Falhas}
    
        A garantia de alta disponibilidade requer que o sistema seja capaz de identificar prontamente quando um componente deixa de operar conforme o esperado.
    
        \subsection{Tipos de Falhas}
        No contexto de sistemas distribuídos industriais, as principais falhas mapeadas incluem falhas de \textit{hardware} (como a queima de componentes ou perda total de energia no servidor), falhas de \textit{software} (\textit{crashes} abruptos do serviço do banco de dados) e problemas de conectividade de rede, que podem variar desde alta latência até o particionamento completo da rede, fenômeno que pode originar cenários de \textit{split-brain}.

        A seguir serão apresentadas duas imagens que demonstram um \textit{cluster} de servidores antes e durante um particionamento de rede:

        \begin{figure}[!ht]
            \conteudoFigura
            [Autor (2026)]
            {1}
            {cenario normal.png}
            {Cluster de Servidores Antes do Particionamento}
            {fig:cenario_normal}
        \end{figure}

        \begin{figure}[!ht]
            \conteudoFigura
            [Autor (2026)]
            {1}
            {cenario anormal.png}
            {Cluster de Servidores Durante um Particionamento} 
            {fig:cenario_anormal}
        \end{figure}
    
        \subsection{Mecanismos de Detecção}
        Para monitorar o estado do sistema, a arquitetura implementa duas camadas de detecção:
        
        \begin{itemize}
            \item \textbf{Monitoramento Local do Banco de Dados}: Através da classe \textit{PostgresHealthChecker}, o agente orquestrador realiza testes de conexão via \textit{Java Database Connectivity} (\verb|JDBC|)\nomenclature{JDBC}{Java Database Connectivity} no banco de dados local a cada três segundos, como ilustrado na Listagem \ref{cod:health-checker}.

\begin{lstlisting}[language=Java, caption={Trecho simplificado da verificação de saúde via JDBC}, label={cod:health-checker}]
private void checkHealth() {
    try (Connection conn = DriverManager.getConnection(url, user, pass);
         PreparedStatement stmt = conn.prepareStatement("SELECT pg_is_in_recovery()");
         ResultSet rs = stmt.executeQuery()) {
        if (rs.next()) {
            boolean isInRecovery = rs.getBoolean(1);
            Status newStatus = isInRecovery ? Status.STANDBY : Status.PRIMARY;
            // Atualiza status local e zera falhas consecutivas
        }
    } catch (SQLException e) { 
        // Incrementa o contador de falhas consecutivas
    }
}
\end{lstlisting}

            Se a conexão falhar por três vezes consecutivas, o banco é classificado como \textit{UNHEALTHY} (não saudável).
            \item \textbf{Monitoramento do Cluster}: A integridade da rede e o status de cada nó são validados através dos \textit{heartbeats} do algoritmo Raft, enviados pelo líder a cada cinco segundos via \verb|gRPC|. A ausência desses sinais indica uma possível falha de rede ou do nó líder, dando início ao processo de eleição.
        \end{itemize}
    
    \section{Caracterização do Ambiente SCADA}
    
        O ambiente \verb|SCADA| que motivou o desenvolvimento desta solução possui restrições e características de alta criticidade. A infraestrutura física da empresa está dividida em dois locais geograficamente distintos, denominados Site A e Site B, cada um abrigando dois servidores físicos e totalizando quatro nós na topologia real de produção.

        A necessidade de manter os dados operacionais absolutamente consistentes é inegociável nesse contexto, pois qualquer perda ou divergência de informação pode impactar diretamente o controle do maquinário industrial. Dado que o protocolo Raft exige um quórum ímpar para garantir a tomada de decisão sem ambiguidade, a topologia de quatro nós apresenta um desafio: em caso de particionamento exatamente ao meio (dois nós de cada lado), o algoritmo seria incapaz de eleger um líder, pois nenhuma das partições atingiria maioria absoluta.

        Para contornar essa limitação sem duplicar o custo de infraestrutura, a arquitetura prevê a inclusão de um \textbf{nó árbitro}, que não armazena dados do PostgreSQL, mas participa das votações do Raft para compor o quórum. Com isso, a topologia real efetiva passa a ter cinco participantes no consenso, mantendo sempre a capacidade de decisão mesmo sob falha de um único nó por localidade. Os testes, a implementação do \textit{raft-ha-manager} e as simulações descritas neste trabalho foram realizados em uma arquitetura reduzida de três nós (dois nós de dados e um árbitro), que reflete fielmente essa mesma lógica de quórum ímpar em escala menor.
    
    \section{Arquitetura de Banco de Dados e Consenso}
    
        A integração entre o banco de dados PostgreSQL e o algoritmo de consenso é a peça central da solução.

        \begin{figure}[!ht]
            \conteudoFigura
            [Autor (2026)]
            {0.35}
            {Novo mural(1).png}
            {Diagrama da Arquitetura do Sistema} 
            {fig:arquitetura}
        \end{figure}
    
        A decisão de desenvolver a máquina de estados do Raft de forma manual, em vez de utilizar uma biblioteca pronta como o \textit{Apache Ratis} ou o \textit{etcd}, foi baseada em três pilares: o aprofundamento acadêmico exigido na pesquisa, a necessidade de integração refinada com o ciclo de vida do PostgreSQL e um maior controle sobre o comportamento de tolerância a falhas. Bibliotecas genéricas dificultariam a implementação de regras de negócio específicas, como bloquear a candidatura de um nó cujo banco de dados está fora do ar.
    
        A comunicação entre os nós do cluster foi desenvolvida sobre o protocolo \verb|gRPC|, utilizando serialização binária com \textit{Protocol Buffers} rodando sobre HTTP/2, baseada em definições ilustradas na Listagem \ref{cod:raft-proto}. 
        
\begin{lstlisting}[language=Java, caption={Definição parcial do serviço Raft em Protocol Buffers (\texttt{.proto})}, label={cod:raft-proto}]
message RequestVoteRequest {
  int64 term = 1;
  string candidate_id = 2;
  int64 last_log_index = 3;
  int64 last_log_term = 4;
}
\end{lstlisting}

        Essa escolha arquitetural provê uma latência de rede significativamente menor do que as tradicionais APIs baseadas em \textit{Representational State Transfer} (\verb|REST|)\nomenclature{REST}{Representational State Transfer} com \textit{JavaScript Object Notation} (\verb|JSON|)\nomenclature{JSON}{JavaScript Object Notation}, característica importante para o tráfego contínuo dos \textit{heartbeats}. O armazenamento do \textit{log} persistente do Raft, no entanto, foi estruturado em formato \textit{JSON Lines} (\verb|.jsonl|), pois permite a adição de novas entradas (\textit{append-only}) de forma rápida, ao mesmo tempo em que preserva a legibilidade para fins de auditoria e depuração.

        Apesar da praticidade e velocidade da adição contínua de registros, o acúmulo indiscriminado de logs ao longo do tempo levaria ao esgotamento do espaço em disco. Para contornar esse problema crítico de infraestrutura e garantir a viabilidade da solução a longo prazo, sistemas baseados no algoritmo Raft implementam rotinas de compactação de log ou \textit{Snapshotting}. Neste processo, o estado atual do sistema é consolidado em um arquivo de captura pontual, também conhecido como \textit{snapshot}, e todos os registros antigos do arquivo \verb|.jsonl| que já foram replicados com segurança são descartados. Isso libera o espaço físico do disco dos servidores sem comprometer o histórico ou a consistência da máquina de estados do cluster.
    
    \section{Dinâmica de Failover}
    
        O processo de recuperação de falhas foi projetado para atuar rapidamente e, acima de tudo, tentar garantir a consistência dos dados, mitigando o problema de \textit{split-brain}.

        O link para acessar o fluxo de failover está disponibilizado em seguida: link do lucidchart

        % https://lucid.app/lucidchart/1bc64d7d-3b8c-4e1b-aa54-d322ccd4938d/edit?invitationId=inv_2b8017e2-f4eb-4128-ae5c-7f2e5129afdf
    
        Quando o monitoramento local detecta que o PostgreSQL do nó líder falhou, ocorre um \textit{Step-Down} voluntário. O nó líder renuncia imediatamente à sua posição, tornando-se um seguidor (\textit{FOLLOWER}) e cessando o envio de \textit{heartbeats}, como demonstrado na Listagem \ref{cod:step-down}.

\begin{lstlisting}[language=Java, caption={Transição de estado do Raft para FOLLOWER (Step-Down)}, label={cod:step-down}]
private void becomeFollower(long newTerm, String leaderId) {
    this.currentTerm = newTerm;
    this.votedFor = null;
    this.currentLeaderId = leaderId;
    this.state = NodeState.FOLLOWER;
    resetElectionTimer();
    stopHeartbeats();
}
\end{lstlisting}

        Esse mecanismo previne que um líder sem banco de dados operacional trave o sistema. Com a interrupção dos \textit{heartbeats}, os \textit{election timeouts} ou temporizadores de eleição, que foram definidos com temporizadores entre quinze e trinta segundos, dos nós restantes expiram, iniciando uma nova eleição. O nó eleito então executa o utilitário \verb|pg_promote|, elevando seu banco de dados local da condição de réplica (\textit{STANDBY}) para primário (\textit{PRIMARY}), lógica demonstrada na Listagem \ref{cod:become-leader}.

\begin{lstlisting}[language=Java, caption={Transição de estado para LÍDER e promoção do banco local}, label={cod:become-leader}]
private void becomeLeader() {
    String previousLeader = currentLeaderId;
    role = RaftRole.LEADER;
    currentLeaderId = config.getNodeId();

    electionTimer.pause();
    startHeartbeatScheduler();

    // Promove o PostgreSQL apenas se o líder anterior tiver falhado
    if (!config.isArbiter()) {
        if (previousLeader != null && !previousLeader.equals(config.getNodeId())) {
            replicationExecutor.submit(this::checkAndExecutePostgresPromotion);
        }
    }
}
\end{lstlisting}
        
        \subsection{Prevenção de Split-Brain}
            Para cenários de particionamento de rede, onde um líder antigo isolado não tem conhecimento de que um novo líder foi eleito, essa situação foi descrita por \citeonline{kleppmann2017designing} como \textit{delusional leader}, o sistema utiliza uma combinação de \textit{Fencing} e \textit{Self-Healing}:
        
            \begin{itemize}
                \item \textbf{\textit{Fencing}}: Se a máquina de estado do Raft de um nó estiver na condição de \textit{FOLLOWER}, mas seu PostgreSQL local relatar estar como \textit{PRIMARY}, o agente orquestrador intervém de forma forçada interrompendo o serviço do banco de dados local (Listagem \ref{cod:fencing}). 

\begin{lstlisting}[language=Java, caption={Execução do Fencing acionando o sistema operacional}, label={cod:fencing}]
public void fenceAndRecover(String leaderIp) {
    log.warn("Executando Fencing: sudo systemctl stop postgresql");
    ProcessBuilder fencePb = new ProcessBuilder(
        "sudo", "/bin/systemctl", "stop", "postgresql");
    Process fenceProc = fencePb.start();
    fenceProc.waitFor(15, TimeUnit.SECONDS);
    // Prossegue com o Auto-Healing...
}
\end{lstlisting}

                Esse isolamento abrupto impede fisicamente que sistemas ou aplicações clientes interajam com um servidor diferente do primário, garantindo a integridade dos dados e seguindo o princípio \textit{Fail-Safe} \cite{kleppmann2017designing};
                
                \item \textbf{\textit{Self-Healing}}: Imediatamente após a parada forçada,  um processo de recuperação automática é acionado. Este script descarta o diretório de dados corrompido ou desatualizado e realiza um clone completo do banco de dados líder legítimo do cluster, reiniciando o serviço como um nó \textit{STANDBY} seguro \cite{toeroe2012service}.
            \end{itemize}

        \subsection{Redirecionamento de Clientes}
        
            A promoção interna de um novo nó ao papel de primário não é suficiente se as aplicações que interagem com os servidores não souberem para qual servidor enviar seus pedidos de gravação, busca, atualização e remoção de dados. Para solucionar a descoberta de serviço de forma prática, sem necessidade de reconfiguração manual nos sistemas, a arquitetura prevê a utilização de um serviço que faz com que as aplicações sigam automaticamente o servidor primário. Quando ocorre um \textit{failover} e um novo líder é estabelecido, mecanismos de rede assumem a responsabilidade de migrar o endereço automaticamente para a nova máquina primária. Com isso, as aplicações SCADA continuam disparando as requisições sempre para o mesmo endereço correto, enquanto o roteamento interno da rede cuida de direcionar o tráfego para o nó ativo naquele instante, diminuindo o tempo total de indisponibilidade dos serviços.
        
    \section{Ambiente de Simulação}
    
        O ecossistema utilizado para testar a aplicação consistiu em uma conjunto de três Máquinas Virtuais (\verb|VM|)\nomenclature{VM}{Virtual Machine} orquestradas no \textit{software} VirtualBox. A automação da criação do ambiente foi realizada por meio do script \verb|vm-create.ps1|, responsável por configurar as máquinas dentro de uma rede \textit{Host-Only}, isolando o ambiente de interferências externas.
    
        O processo de implantação e configuração do PostgreSQL na versão 16, e do \textit{raft-ha-manager}, executado em uma \textit{Java Virtual Machine} (\verb|JVM|)\nomenclature{JVM}{Java Virtual Machine} 21, foi automatizado através de \textit{shell scripts} executados de forma sequencial:
        
        \begin{enumerate}
            \item \verb|01-base.sh|: realiza a configuração inicial do sistema operacional em todos os nós, definindo o endereço IP estático, ajustando as regras de \textit{firewall}, editando o arquivo \verb|/etc/hosts| para a resolução de nomes locais, além de instalar o Java e criar o usuário de serviço;
            
            \item \verb|02-postgres-primary.sh|: configura o PostgreSQL como primário no primeiro nó, ajustando os parâmetros de rede no arquivo \verb|postgresql.conf| e as regras de acesso no arquivo \verb|pg_hba.conf|, finalizando com a criação dos usuários de replicação;
            
            \item \verb|03-postgres-standby.sh|: atua nos demais nós de dados para configurá-los como réplicas, realizando a clonagem inicial através da ferramenta \verb|pg_basebackup| e estabelecendo o canal de \textit{streaming replication};
            
            \item \verb|04-deploy-raft.sh|: implanta o executável Java da aplicação, distribui os arquivos de configuração no formato \verb|YAML| e registra o \textit{raft-ha-manager} como um serviço do \textit{systemd}, garantindo a sua execução contínua em \textit{background}.
        \end{enumerate}
        
        Essa padronização garantiu que os testes de \textit{failover}, \textit{fencing} e \textit{self-healing} fossem validados em um ambiente controlado e replicável, espelhando os desafios operacionais do ambiente \verb|SCADA| industrial.

        \subsection{Limitações da Simulação}

        A literatura acadêmica clássica do algoritmo Raft, concebida por \citeonline{raft2014}, determina que os \textit{heartbeats} devem ser enviados em intervalos de 10 a 50 milissegundos com \textit{timeouts} de eleição oscilando entre 150 e 300 milissegundos. Tais valores garantem um tempo de inatividade da rede quase imperceptível, requisito importante para um ambiente industrial.

        Contudo, exclusivamente para a prova de conceito e validação funcional neste trabalho acadêmico, os valores foram configurados de forma conservadora: \textit{heartbeats} de 5 segundos e temporizadores de eleição entre 15 e 30 segundos. Estes ajustes nos tempos foram necessários para compensar a limitação introduzida pelo próprio ambiente de simulação. Como os testes foram realizados em \verb|VM|s compartilhando recursos do mesmo processador físico, somando isso a execução do orquestrador na \verb|JVM|, o sistema tornou-se suscetível a pequenas pausas de escalonamento de processos do hospedeiro e a rotinas de coleta de lixo, processo que ocorre automaticamente para limpar memória nas \verb|JVM|s. Utilizar os valores de milissegundos no ambiente virtualizado geraria latências falsas e expirações prematuras, provocando eleições de líder desnecessárias e instabilidade no \textit{cluster}. Dessa forma, a adoção de valores maiores permitiu validar com maior facilidade e segurança a lógica de \textit{failover} sem essas pequenas interferências da virtualização. 